\section{The first full ambient historical normal anchor}

\begin{definition}[Original normal contact and a genuine first step]
\label{hir:disjoint-first-step-normal-history}
\uses{hir:disjoint-coherent-normal-record,hir:disjoint-coherent-normal-initial-value}
\lean{TranscendentalNakamaye.Pending.HironakaDisjointFirstStepHistoricalAnchor.normalContact}
\leanok
Let $n>0$, let $M=(\mathbf C^n,\mathbf C^n,\varnothing,I,1)$ lie in the
depth-$q$ invariant domain, and put $J=J(I)$ in the whole ambient space
$\mathbf C^{n+1}$. Let $C$ be empty or the first maximum locus of the
depth-$q$ invariant of $M$. Choose an actual blowup $\Pi$ along $s(C)$ and
write $S$ for this one-step tower from the original stabilized space. Put
$O=\mathbf C^{n+1}$, regarded as the whole original open, and let $K_t$ be
the original contact marked ideal for the literal coordinate $t$, formed
before this blowup and restricted to $O$. In every value below the starting
marked ideals are these original $J$ and $K_t$; the tail is the genuine $S$.
\end{definition}

\begin{lemma}[No divisorial support gives unit monomial stalks]
\label{hir:disjoint-full-ambient-unit-monomial-stalk}
\lean{TranscendentalNakamaye.Pending.HironakaDisjointFirstStepHistoricalAnchor.one_mem_monomial_of_empty_codimOne}
\leanok
Let $A$ be a coherent ideal on a complex manifold of dimension $m+1$,
with whole carrier and mark one, of maximal order one. Suppose its
codimension-one cosupport part is empty. For any simple normal crossings
boundary $E$, the constant one belongs to the monomial-part stalk at every point.
\end{lemma}
\begin{proof}
Fix a point $y$ and a member $F$ of $E$ through $y$. Suppose its exponent
is positive. Divisibility by its actual exponent, followed by lowering the
exponent to one, gives a neighborhood of $y$ on which $F$ is contained in
the cosupport. Choose an open neighborhood $V$ with this containment.
Simple normal crossings makes $F$ smooth of dimension $m$ at $y$, hence
$F\cap V$ is smooth of dimension $m$ there. Its containment in the analytic
cosupport gives local dimension at least $m$ for that cosupport. The whole
carrier has dimension $m+1$ and the marked ideal has maximal order one;
therefore the codimension-one criterion puts $y$ in its codimension-one
part, a contradiction. Thus every boundary member through $y$ has exponent
zero. Its zeroth ideal power contains one. These memberships give one in
the monomial-part stalk by the definition, with no condition from members missing $y$.
\end{proof}

\begin{lemma}[The actual upper first stage and its original clocks]
\label{hir:disjoint-first-step-upper-clocks}
\uses{hir:disjoint-first-step-normal-history,hir:disjoint-full-ambient-unit-monomial-stalk,hir:disjoint-coherent-normal-initial-value}
\lean{TranscendentalNakamaye.Pending.HironakaDisjointFirstStepHistoricalAnchor.firstStep_upper_isCanonical,TranscendentalNakamaye.Pending.HironakaDisjointFirstStepHistoricalAnchor.original_members_nil,TranscendentalNakamaye.Pending.HironakaDisjointFirstStepHistoricalAnchor.original_level_zero,TranscendentalNakamaye.Pending.HironakaDisjointFirstStepHistoricalAnchor.firstStep_current_boundary,TranscendentalNakamaye.Pending.HironakaDisjointFirstStepHistoricalAnchor.firstStep_nmOrd_eq_one,TranscendentalNakamaye.Pending.HironakaDisjointFirstStepHistoricalAnchor.firstStep_not_special,TranscendentalNakamaye.Pending.HironakaDisjointFirstStepHistoricalAnchor.firstStep_order_run,TranscendentalNakamaye.Pending.HironakaDisjointFirstStepHistoricalAnchor.firstStep_step_run}
\leanok
The actual $S$ is canonical for the depth-$(q+1)$ invariant of $J$.
Its current boundary is exactly $[\Pi^{-1}s(C)]$. The original-member list
is empty. At every point of the transformed cosupport, its nonmonomial
order is one, its original level is zero, and it is nonspecial. Both its
original order run and its original Step One run have length one, so their
heads are the original identity tower and their tails are the whole $S$.
\end{lemma}
\begin{proof}
The initial normal recoding identifies the full ambient first maximum
with $s(C)$ when $C$ is the base first maximum. If $C$ is empty, $s(C)$ is
empty. Thus in either allowed case the actual first blowup is canonical.
The domain theorem makes it an actual test tower. The transformed boundary
is the strict transforms of the empty original list followed by the
inverse-image centre, proving the asserted singleton including empty centres.
Original members are the prefix of length zero of this current list, hence
the original level is zero.

The initial codimension-one part is empty because the initial stabilized
record is nonspecial at every point. Maximal order is one, so its
codimension-one part remains empty along this actual test tower. The same
fact rules out the level-zero special branch, while the positive-level
special branch contradicts level zero. Maximal order also persists along
the test. Apply the preceding unit-monomial-stalk lemma to the current
coherent ideal and its actual simple normal crossings boundary. For a member
$g$ of its nonmonomial part at a cosupport point, multiply by the monomial
member one in the defining colon. This puts $g$ in the original ideal on
the whole carrier, so $g$ vanishes at that point. Thus every member of the
nonmonomial part vanishes, giving nonmonomial order at least one. Canonical order
monotonicity bounds it by the nonmonomial order of its original image,
which is one by the initial normal order theorem. It is therefore one.
The original image is in the original cosupport, where both the original
order-one regime and nonspecial level-zero regime hold. The recursive
definition of the run for this real one-step tower adds one to its zero
initial run. The recursive split then has original head and actual full tail.
\end{proof}

\begin{lemma}[The literal original contact has the same real first centre]
\label{hir:disjoint-first-step-literal-contact-canonical}
\uses{hir:disjoint-first-step-normal-history,hir:disjoint-coherent-normal-whole-trace,hir:disjoint-coherent-normal-initial-value}
\lean{TranscendentalNakamaye.Pending.HironakaDisjointFirstStepHistoricalAnchor.firstStep_contact_isCanonical,TranscendentalNakamaye.Pending.HironakaDisjointFirstStepHistoricalAnchor.firstStep_contact_cosupport,TranscendentalNakamaye.Pending.HironakaDisjointFirstStepHistoricalAnchor.firstStep_contact_boundary}
\leanok
The original $K_t$ lies in the depth-$q$ domain and its restricted actual
one-step tower is canonical. Its transformed cosupport is exactly the
whole-open trace of the transformed cosupport of $J$. Its current boundary
retains the same born member, with the literal chronological label zero.
\end{lemma}
\begin{proof}
For mark one, homogenization and coefficient formation leave the whole
coherent ideal unchanged. Thus the original literal $t$-contact is precisely
the whole-open restriction of $(H,\varnothing,J(I),1)$. Its natural contact
domain follows from the original full ambient domain, the actual tangent
direction, and absence of the codimension-one part. At the identity tower
both contact tails are canonical identities. The zero-level contact anchor
therefore identifies its finite value with the second component of the
established original normal Step One initial formula. That component is
the entire depth-$q$ base value, a consequence of the whole functional section
identity $s^*J(I)=I$ and the initial section value theorem. The contact-tail
cosupport dictionary at the identity gives exactly the section image of the
base cosupport. Hence its entire initial maximum is $s(C)$ when $C$ is the base
maximum. The empty-centre case remains empty on restriction. This proves
canonicality of its actual first step independently of any future value
comparison. The contact-tail cosupport theorem applies to the literal
original tangent direction and the genuine upper test tower. It gives the
whole asserted transformed cosupport equality. Its transform appends the
inverse-image centre after the empty original list, retaining that born slot.
\end{proof}

\begin{theorem}[Actual Step One value through the original normal contact]
\label{hir:disjoint-first-step-step-one-anchor}
\uses{hir:disjoint-first-step-upper-clocks,hir:disjoint-first-step-literal-contact-canonical}
\lean{TranscendentalNakamaye.Pending.HironakaDisjointFirstStepHistoricalAnchor.firstStep_stepOne_value}
\leanok
At every point $y$ of the full ambient transformed cosupport, Step One
evaluated from the original $J$ along $S$ equals
\[
 \left(0,\;R_q^+(K_t,S|_O,y)\right),
\]
where the second entry is finite and its lower value is taken along the
actual one-step original contact history.
\end{theorem}
\begin{proof}
At the original stage, the established normal Step One initial value is
$(0,R_q(M,\mathrm{id},x))$. Thus its first maximum also equals the original
base first maximum lifted by $s$. This makes the same real $S$ canonical
for the original maximal-order Step One function. The canonical-tail
theorem supplies the actual selected lower record's domain, canonical tail,
mask membership, and cosupport membership. The upper-clock lemma identifies
the selected split as original head and real one-step tail, with level zero.
The literal coordinate $t$ is an actual original tangent direction on the
whole ambient open; the empty original-member list misses that open. The
preceding lemma supplies the independently proved domain and canonicality
of its literal contact tail. Therefore the zero-level contact anchor
identifies the selected lower value with the actual value for $K_t$ along
$S|_O$. Nonspecialness selects the finite branch in the definition of Step
One. Substitution gives the displayed equality.
\end{proof}

\begin{theorem}[The first full ambient historical anchor]
\label{hir:disjoint-first-step-full-historical-anchor}
\uses{hir:disjoint-first-step-upper-clocks,hir:disjoint-first-step-step-one-anchor}
\lean{TranscendentalNakamaye.Pending.HironakaDisjointFirstStepHistoricalAnchor.firstStep_historical_value,TranscendentalNakamaye.Pending.HironakaDisjointFirstStepHistoricalAnchor.firstStep_historical_maximum_locus}
\leanok
The original full ambient invariant after this actual first blowup satisfies
\[
 R^+_{q+1}(J,S,y)=\iota_q\bigl(R_q^+(K_t,S|_O,y)\bigr)
\]
on its entire transformed cosupport. Its next maximum locus is precisely
the image of the maximum locus of this literal contact historical value.
This assertion includes empty transformed cosupports.
\end{theorem}
\begin{proof}
The upper-clock lemma puts the general invariant in its positive-order
branch with order one, with companion formed at the original stage and
the whole real $S$ as its tail. The initial order-one mask is the whole
original ambient open. On it the companion of the initial empty-boundary
mark-one record is restricted-equivalent to the original $J$: its
nonmonomial ideal with empty boundary is the restriction of the whole ideal,
and the companion uses the branch where order equals the mark. Its domain
and maximal order follow from the companion-domain theorem. The original
Step One tail is canonical by the preceding initial-centre argument;
restriction to the whole order mask and restricted-equivalence invariance
therefore make the companion tail canonical. Representative invariance and whole-open
restriction invariance identify its value with original Step One for $J$
along $S$. The preceding theorem gives $(0,R_q^+(K_t,S|_O,y))$ for this value.
Substitution into the general positive-order constructor gives exactly
$\iota_q$ of that actual contact historical value, with both finite
inclusions. The contact cosupport dictionary identifies all points on which
these values are compared. Since $\iota_q$ is an order embedding, a point
dominates that entire contact cosupport exactly when its image dominates
the entire full ambient transformed cosupport. This proves the maximum
locus equality. If either cosupport is empty, the dictionary makes both
empty and both defining maximum loci are empty.
\end{proof}

\begin{remark}[The boundary correction in a further comparison]
This historical anchor does not identify its contact value with the raw
base value after the blowup. For a nonbottom monomial branch, the ambient
dimension in its key increases by one while its chronological subset stays
fixed. A further cross-ambient historical identification therefore uses
the recursive same-index dimension map, not the unmodified base value.
For example, the first transform of $(x^2,1)$ on a line has ideal $(x)$ and
born boundary $[x=0]$, with monomial key $(0,1,\{0\})$. The first $x$-chart
transform of $(x^2,t)$ has contact ideal $(x)$ in a plane and key
$(1,1,\{0\})$. The theorem retains the actual lower history and does not
restart the current marked ideal with that born member as original boundary.
\end{remark}
